%% ============================================================
%%  §7 结论
%%  要证明的一句话（全篇落点）：
%%    能力不等于效果：MLIR 让层可协商，但收益的门槛恰好落在现实 AI 框架的数量级附近，
%%    这既解释了它近年被广泛采纳，也解释了为什么采纳不是自动发生的。
%%  建议篇幅：0.4 页
%% ============================================================

\section{结论}
\label{sec:conclusion}

%% 三段结构：
%% ① 回到命题：一句话概括本文把 N×M 叙事重构为什么。
%% ② 用数字收口：|E| = 76（下界）；收敛在端点（llvm 104 / 19 源）；门槛 4.6（架构级）
%%    到 15.5（工程级，p=0.6），敏感性 2.9 倍（8.7→25.6）。
%% ③ 落点：能力 != 效果；治理成本（§5）是 p 上不去的机制原因；
%%    由此可解释"MLIR 在 2020–2023 被广泛采纳，而更早的单层 IR 方案没有"。

%% 可选一句展望：把 torch-mlir 与 OpenXLA 的语料并入，可把"栈顶碎片化"从间接推断
%% 变成直接测量。
